\section*{How to Use the Science Without Blaming Yourself}\label{sec:note_for_readers_text}
\addcontentsline{toc}{section}{How to Use the Science Without Blaming Yourself}
\noindent Every chapter you just read can feel like an accusation when you are already struggling to fall asleep, work shifts,
keep children alive, or manage ADHD with stimulants. The point of this book is to show that sleep is a phase of thinking, not a
moral test. That has three consequences for readers who live inside chronic exhaustion:

\begin{enumerate}
    \item \textbf{Stimulants and late-night hyperfocus are not proof of weakness.} They are attempts to prop up a brain whose
    offline maintenance cycle keeps getting deferred. The evidence in Chapter~\ref{ch:what_fails} shows that deprived brains
    preserve low-level skills but lose meta-control; the frustrating ``why can't I just go to bed?'' loop is neurobiology, not
    character.
    \item \textbf{Track obstacles, not just hours.} Sleep diaries matter less for their totals than for the patterns they reveal:
    shift rotations, caretaking emergencies, ambient noise, or the rebound insomnia that follows discontinuing medication. Name
    the constraints so that any intervention you try---CBT-I, light therapy, or simply negotiating for protected time---is aimed
    at the right target.
    \item \textbf{Treat small, repeatable rituals as calibration, not hacks.} When recovery sleep finally arrives, note what made
    it possible (a darkened room, a trusted partner taking the next feeding, stepping away from screens an hour early). Those are
    not cures, just evidence that your nervous system still knows how to cycle when it is given a chance.
\end{enumerate}

The rest of the book returns to experimental detail, but keep this note handy. You are not being graded on compliance; you are
learning how to protect a biological algorithm that civilization keeps interrupting.
